%%%BLOCK id=071-01 ch=02 sec=重力弹力摩擦力 kind=summary alt=none cont=none y=0.09-0.40
% 页面左侧边缘露出的「水/飞/火/汽车/水平力的合成与分解/圆/水平/单摆」等字属相邻页边缘，未转录
\textbf{（页边）相互作用}

\textbf{重力\quad 弹力\quad 摩擦力}
\begin{itemize}
\item 弹力的分析与计算
  \begin{itemize}
  \item 用胡克定律
  \item 力的平衡条件
  \item 牛顿第二定律
  \end{itemize}
\item 摩擦力的有无与方向判断：假设法\quad 状态法\quad 转换法
\item 摩擦力大小的计算
  \begin{itemize}
  \item 滑摩：$f_{\text{滑}}=\mu N$\quad $\pm$\quad $\mu mg$ % CHECK: 「=μN」与「μmg」之间的符号像「±」（也可能是「=」），按原样转录
    \\ 或由平衡条件、牛二求
  \item 静摩：一般 $f_{\text{静max}}=\mu N$\quad 有时 $f_{\text{静max}}>\mu N$
    \\ 或由平衡条件、牛二求.
  \end{itemize}
\item 摩擦力突变问题：快带慢\quad 慢阻快\quad $+$\quad 快后慢\unclear{断}
  \\ $\leftarrow$ $v$的关系变了\quad （手写箭头由「$v$的关系变了」指向「摩擦力突变问题」）
\end{itemize}
%%%END
%%%BLOCK id=071-02 ch=02 sec=力的合成与分解 kind=summary alt=none cont=none y=0.43-0.60
\textbf{力的合成与分解.}\quad \textbf{共点力}

\begin{tabular}{@{}lll@{}}
力的合成 & 等效替代法 & 平行四边形法则 \\
 & 平移 & 三角形法则 \\
力的分解 & 正交分解. & \\
\end{tabular}
%%%END
%%%BLOCK id=071-03 ch=02 sec=活结死结与动杆定杆 kind=knowledge alt=none cont=none y=0.63-0.71
\textbf{活结与死结、动杆与定杆模型}

活结两侧绳张力相等、死结绳两侧力不一定相等、动杆弹力必沿杆，定杆弹力方向任意
%%%END
%%%BLOCK id=071-04 ch=02 sec=受力分析与共点力平衡 kind=summary alt=none cont=next y=0.71-1.00
\textbf{受力分析、共点力的平衡}
\begin{itemize}
\item \unclear{顺序}：重弹摩场、合成法\quad 正交分解法.
\item 静态平衡
\item 临界：极限分析法、数学分析法（二次函数、三角函数、不等式）\quad 图解法.
\item 动态平衡.
  \begin{itemize}
  \item 解析法
  \item 图解法
  \item 相似三角形法.
  \item 正交分解法
  \item 力辅助圆 % 页面底部被照片截断，此行只拍到一半
  \end{itemize}
\end{itemize}
%%%END
%%%BLOCK id=072-01 ch=02 sec=受力分析与共点力平衡 kind=summary alt=none cont=prev y=0.02-0.60
% 本页为横置页（已转正）；左缘倾斜的小字、红色封面等属相邻页/杂物，未转录
\textbf{力学\quad 共点力平衡}

必然\quad 所有力延长线交于一点

\medskip
\textbf{静态平衡}\quad （求某个力的值）
% 手写箭头：「在一条线上」右侧有箭头指向「等大反向」；另有一斜细箭头自此处指向下方「三力」一行；「力三角形闭合」「力的多边形闭合」与「建系…」之间有一个空心粗箭头
\begin{itemize}
\item 二力\quad 在一条线上 $\to$ 等大反向
\item 三力\quad 力三角形闭合 $\Rightarrow$ 建系\ $\sum F_x=0$\ \ $\sum F_y=0$
\item 多力\quad 力的多边形闭合 $\Rightarrow$ 建系\ $\sum F_x=0$\ \ $\sum F_y=0$
\item （页面右侧另有一列，自上而下）有直角用勾股定理 / 非直角三角形 / 正弦定理 / 余弦定理
\end{itemize}

\medskip
\textbf{动态平衡}\quad （求力的变化）
\begin{itemize}
\item 二力\quad 等大反向
\item 三力\quad 力三角形\quad 受力分析：
  \begin{itemize}
  \item 旋转\quad 2个力方向不变\quad 、延长，
  \item 相似\quad 有绳子/杆\quad 非特殊三角形\quad 相似\ 正\unclear{余}弦定理
  \item 力辅助圆\quad 有一角不变\quad 有\unclear{个}力不变
  \item 建系分解.
  \end{itemize}
\item 多力 % 位置在「动态平衡」之下、与二力三力同列
  \begin{itemize}
  \item 4个力\quad 有摩擦\quad 有支持力\quad 用摩擦角
  \item 没有摩擦角\quad 正交分解到$X$、$Y$方向\,方程
  \end{itemize}
\end{itemize}
%%%END
%%%BLOCK id=072-02 ch=02 sec=受力分析与共点力平衡 kind=method alt=none cont=none y=0.42-0.65
\textbf{力学两确定.}
\begin{itemize}
\item 平衡还是非平衡
\item 研究对象是谁
\end{itemize}
%%%END
%%%BLOCK id=086-02 ch=02 sec=共点力平衡·临界极值 kind=problem alt=03 cont=none y=0.365-0.45
\begin{problem}[1、]
%FIG p086_d 0.085 0.368 0.245 0.442
\notefig[0.3]{figures/p086_d.png}
\unclear{$\mu$}$=0.125$\quad 系统静止，求 $m_B$ 范围
% CHECK: 手写像 "M=0.125"，按物理应为动摩擦因数 μ=0.125（斜面 37°，A、B 由绳经斜面顶端定滑轮相连，A 在 B 上）；答案中的 M 即 m_B，m 为 m_A
\begin{solution}
\[ \begin{cases} m_{B\,\max}\text{ 时 }f_B\text{ 向上}\\[2pt] m_{B\,\min}\text{ 时 }f_B\text{ 向下}\end{cases}\qquad \frac37m\le M\le\frac95m \]
\end{solution}
\end{problem}
%%%END
%%%BLOCK id=114-03 ch=02 sec=共点力平衡 kind=problem alt=90 cont=none y=0.44-0.58
\begin{problem}[3]
% 图p114_d内含标注：水平线上的T、θ、mg、F方向等（受力示意）；图p114_e：F、θ、mg、T 的力三角形
%FIG p114_d 0.04 0.44 0.262 0.54
\notefig[0.35]{figures/p114_d.png}

$F=\rho SV^2\sin^2\theta$\quad 求 $V_{\min}$
\begin{solution}
\[mg=F\cos\theta\]
\[\therefore\ V=\sqrt{\frac{mg}{\rho S\sin^2\theta\cos\theta}}\qquad \tan\theta=\frac{\sqrt{2}}{2}\ \text{时}\ V_{\min}\] % CHECK: 若 F∝sin^2θ，使 sin^2θ cosθ 最大应得 tanθ=√2；原稿写作 √2/2（拍照看似 √2 over 2）

%FIG p114_e 0.70 0.495 0.84 0.585
\notefig[0.25]{figures/p114_e.png}
\end{solution}
\end{problem}
%%%END
%%%BLOCK id=163-03 ch=02 sec=受力分析·接触面作用力 kind=note alt=none cont=none y=0.33-0.39
%FIG p163_c 0.57 0.33 0.708 0.39
\notefig[0.25]{figures/p163_c.png}
$M$为$m$\ 作用力\quad $F=\sqrt{N^2+f^2}$ % CHECK: “为”字疑为“对”（M对m作用力）；m 后有涂改液痕迹
%%%END
%%%BLOCK id=171-01 ch=02 sec=验证力的平行四边形定则·实验 kind=method exp=1 alt=none cont=none y=0.02-0.22
\textbf{力的平行四边形}

①

%FIG p171_a 0.14 0.043 0.47 0.185
\notefig[0.4]{figures/p171_a.png}

记录$F$大小与方向

$F_3$合力：橡皮/测力计
%%%END
%%%BLOCK id=171-02 ch=02 sec=验证力的平行四边形定则·实验 kind=method exp=1 alt=none cont=none y=0.23-0.43
② 橡皮筋

%FIG p171_b 0.235 0.258 0.49 0.402
\notefig[0.35]{figures/p171_b.png}

在力的方向上画伸长量

%FIG p171_c 0.49 0.35 0.63 0.41
\notefig[0.2]{figures/p171_c.png}

\uwave{不用让$O$在同一位置了}
%%%END
%%%BLOCK id=171-03 ch=02 sec=验证力的平行四边形定则·实验 kind=method exp=1 alt=none cont=none y=0.43-0.68
③ 三力 $\to$ 转$+$相似$+$做圆

%FIG p171_d 0.14 0.44 0.51 0.64
\notefig[0.4]{figures/p171_d.png}

% 红笔批改：A 行末有红勾；B 行末红×并补写「和方向」；C 行末有红勾/红线并旁注「用细绳标方向」；D「大一些」下有红线+红勾
A：$O$点不能动；同次\ \red{$\checkmark$}

B：记下测力计读数，$O$点位置\red{$\times$和方向} % 「记下」为行间小字插入

C：拉橡皮筋的细绳长一些，效果好\ \red{$\checkmark$}\quad \red{用细绳标方向}

D：拉力要适当大一些\ \red{$\checkmark$}
%%%END
%%%BLOCK id=171-04 ch=02 sec=弹簧·胡克定律实验 kind=note exp=1 alt=none cont=none y=0.68-0.73
结点$D$偏高没有影响，直$R$斜\unclear{放}$k$偏小 % CHECK: 「D」疑为「O」；「直R」疑为「直尺」；该句疑指弹簧(k)实验中直尺斜放读数偏大 → k偏小
%%%END
%%%BLOCK id=178-01 ch=02 sec=胡克定律·实验 kind=method exp=1 alt=none cont=none y=0.05-0.34
\textbf{胡克定律}

但凡验证必用图象(图象\unclear{，}必为直线)

\[
F=k\cdot\Delta X\qquad
\left\{
\begin{aligned}
F_1&=kX_1\to\text{形}\\
F_2&=kX_2\to\text{形}
\end{aligned}
\right.
\]

\[
\Delta F=k\cdot\Delta X=k\left[X_1'-X_0-(X_2'-X_0)\right]=k\Delta X
\]

$\downarrow$ 现在$K$\unclear{…} % CHECK: 此行有箭头自 kΔX 指下；字迹被涂黑，只辨出「现在K」

2011\ 拉$F_1,X_1$\ 压$F_2,X_2$\ 求$k$\qquad $k=\dfrac{\overset{\text{力换方向}}{F_1-(-F_2)}}{X_1-X_2}$

操作\quad $F_1\Rightarrow\overset{\text{形变}}{\Delta X_1}$\qquad $k_1=\dfrac{F_1}{\Delta X_1}$

法一：\ $F_2\to\Delta X_2$\qquad $k_2=\dfrac{F_2}{\Delta X_2}$\qquad 求平均$\overline{k}=\dfrac{k_1+k_2+\cdots+k_n}{n}$

$\vdots$
%%%END
%%%BLOCK id=178-02 ch=02 sec=胡克定律·实验 kind=method exp=1 alt=none cont=none y=0.34-0.72
法二：图象
\begin{keybox}
\[F=k(X-X_0)\]
\end{keybox}

%FIG p178_b 0.37 0.335 0.80 0.495
\notefig[0.43]{figures/p178_b.png}

\begin{keybox}
\[F=k(X_0-X)\]
\end{keybox}

%FIG p178_c 0.40 0.497 0.76 0.65
\notefig[0.36]{figures/p178_c.png}

\begin{center}
\begin{tabular}{c|c|c|}
$F$ & $F_1$ & $F_2$ \\
\hline
$X$ & $X_1$ & $X_2$ \\
\end{tabular}
\end{center}
%%%END
%%%BLOCK id=178-03 ch=02 sec=胡克定律·实验 kind=method exp=1 alt=none cont=none y=0.72-0.83
法三：逐差

\[
\Delta\overline{X}=\frac{X_7+X_6+X_5+X_4-(X_3+X_2+X_1+X_0)}{16}\quad;\quad F=k\Delta\overline{X}\quad k=\frac{F}{\Delta\overline{X}}
\]
%%%END
%%%BLOCK id=178-04 ch=02 sec=弹簧串并联 kind=formula alt=none cont=none y=0.82-0.94
\begin{keybox}
串联\quad $\dfrac{1}{k'}=\dfrac{1}{k_1}+\dfrac{1}{k_2}$\qquad 串$n$个\quad $k'=\dfrac{k_0}{n}$

并联\quad $k'=nk_0$\quad 并$n$个
\end{keybox}
%%%END
%%%BLOCK id=179-05 ch=02 sec=测动摩擦因数·实验 kind=method exp=1 alt=03 cont=none y=0.51-0.60
\fbox{测$\mu$}

%FIG p179_g 0.01 0.536 0.235 0.598
\notefig[0.25]{figures/p179_g.png}

能测出$A$与$B$间$\mu=\dfrac{F}{m}$ % CHECK: 手写为 F/m（疑漏 g，应为 F/mg）
%%%END
%%%BLOCK id=181-05 ch=02 sec=弹力·测杨氏模量 kind=knowledge exp=1 alt=18 cont=none y=0.69-0.78
\textbf{测杨氏模量}\quad $F=kx$\quad $k=Y\dfrac{S}{L}$\quad $Y$单位是 $\mathrm{Pa}$。

$L$为未受拉力时长度，$S$为横截面积。$\leftarrow$ 测直径 $D$ 用螺旋测微器 % “D”写在“测直径”上方行间

\[ Y=\frac{kL}{\pi\dfrac{D^2}{4}} \]
%%%END
%%%BLOCK id=184-07 ch=02 sec=验证力的平行四边形定则 kind=method exp=1 alt=none cont=none y=0.82-1.00
\textbf{验\unclear{证力的}平行四边形定则} % “证力的”为标题上方小字插入

%FIG p184_g 0.045 0.87 0.20 0.998
\notefig[0.25]{figures/p184_g.png}
% 图中标注：F_1、F_2、F、F_3，底部写有“F理 F真(测?)”

比较 $F$ 与 $F_3$ 的大小与方向，得出结论：在误差允许范围内，力的平行四边形\unclear{法}则成立
%%%END
%%%BLOCK id=186-01 ch=02 sec=验证力的平行四边形定则 kind=note alt=none cont=none y=0.03-0.21 exp=1
%FIG p186_a 0.04 0.027 0.22 0.165 soft
\notefig[0.25]{figures/p186_a.png}

要记录结点 $O$ 的位置、每条绳上所挂钩码的个数，$OA$、$OB$、$OC$\unclear{…}（行间小字：大小方向）
% 此行右侧为淡铅笔字迹、对焦模糊：
（右侧淡字：\unclear{分力的大小、方向作用……}；\unclear{$F_B$ 沿绳的方向}）

不用让 $OA\perp OB$，不用让 $O$ 点不变，可以换成弹性绳

（只需稳定即可）\quad（分力和合力可在同一个图示中表示出来）\quad（不影响受力的方向和\unclear{钩码}个数）
%%%END
%%%BLOCK id=221-02 ch=02 sec=受力分析·平衡解题方法 kind=method alt=none cont=none y=0.28-0.38
\textbf{2、}
\begin{itemize}
\item 正交分解
\item 4力\quad 先合成2定力\quad 后三力汇交 \quad 等效替代
\item 定角\quad 画力\quad 辅助圆
\end{itemize}
%%%END
%%%BLOCK id=221-07 ch=02 sec=弹力·图像 kind=method alt=90 cont=none y=0.52-0.63
\textbf{6、}图象分段、曲线、轴点线斜截面
%FIG p221_d 0.055 0.55 0.20 0.62
\notefig[0.25]{figures/p221_d.png}
两个劲度系数（描点时分段了）。
%%%END
%%%BLOCK id=235-04 ch=02 sec=重力·弹力 kind=knowledge alt=03 cont=none y=0.535-0.67
\textbf{重力、弹力}\par
万有引力指向地心，重力竖直向下（无条件）\par
太空完全失重，不能用弹簧测力计测$G$，可以测拉力．天平不能用 \unclear{摆钟}、压强计\par
\unclear{接触}因弹性形变而产生的力\par
力学两确定：\begin{tabular}[t]{@{}l@{}}研究对象\\运动状态\end{tabular}
%FIG p235_h 0.05 0.588 0.225 0.668
\notefig[0.2]{figures/p235_h.png}
无$T$无$N$
%%%END
%%%BLOCK id=236-01 ch=02 sec=滑动摩擦力 kind=knowledge alt=none cont=none y=0.05-0.33
\textbf{滑动摩擦}　产生条件有弹力　$\mu N$

阻碍相对运动

方向与相对运动方向相反，与运动方向不一定相同

$f_{\text{滑}}$不一定是阻力

%FIG p236_a 0.145 0.158 0.275 0.225
\notefig[0.25]{figures/p236_a.png}

%FIG p236_b 0.21 0.208 0.615 0.325
\notefig[0.4]{figures/p236_b.png}
%%%END
%%%BLOCK id=236-02 ch=02 sec=静摩擦力 kind=knowledge alt=none cont=none y=0.34-0.54
\textbf{静摩擦}　$0<f\le f_{\max}$ % CHECK: 手写左端像 "<"（物理上应为 $\le$），照录

$\left\{\begin{array}{l}N\\ \mu\end{array}\right.$

阻碍　相对运动趋势

与运动方向无关

与相对运动趋势方向相反

%FIG p236_c 0.372 0.36 0.68 0.424
\notefig[0.3]{figures/p236_c.png}

\[
\left\{\begin{array}{ll}
F>mg\sin\alpha & \text{下}\\
F=mg\sin\alpha & \text{无}\\
F<mg\sin\alpha & \text{上}
\end{array}\right.
\]

$f_{\text{静}}$方向：
\[
a\left\{\begin{array}{ll}
>g\sin\theta & \text{沿下}\\
=g\sin\theta & \text{无}\\
<g\sin\theta & \text{沿上}
\end{array}\right.
\]
%%%END
%%%BLOCK id=236-03 ch=02 sec=摩擦力突变 kind=method alt=none cont=none y=0.55-0.68
%FIG p236_d 0.10 0.558 0.2445 0.622
\notefig[0.25]{figures/p236_d.png}

$f_{\text{静}}=mg\sin\omega t$　　$f_{\text{滑}}=\mu mg\cos\omega t$

%FIG p236_e 0.225 0.612 0.44 0.675
\notefig[0.25]{figures/p236_e.png}
%%%END
%%%BLOCK id=236-04 ch=02 sec=摩擦力突变 kind=method alt=03 cont=none y=0.69-0.88
% 图中标注：F=kt，μ，f，G，f=μkt，向下减速，静止，向下加速，t
%FIG p236_f 0.045 0.695 0.34 0.875
\notefig[0.3]{figures/p236_f.png}
%%%END
%%%BLOCK id=236-05 ch=02 sec=摩擦力突变 kind=problem alt=03 cont=none y=0.88-1.0
%FIG p236_g 0.11 0.885 0.265 0.995
\notefig[0.25]{figures/p236_g.png}

$F_{\text{拉}}=5\unit{N}$　　$m=5\unit{kg}$　均静止　右加速 $a$ 从 $0\to 2\unit{m/s^2}$ % CHECK: "5" 与 F拉=5N 的手写数字同形，由 $a_1=F/m=1\unit{m/s^2}$ 判断 m=5kg；"均"为插入字

$F+f_{\max}=ma_2$　　$a_2\ge 2\unit{m/s^2}$　　$\therefore F_{\text{拉}}$不变，$f$先向左$\downarrow$ 再向右$\uparrow$

$f=F=5\unit{N}\le f_{\max}$　　$a_1=1\unit{m/s^2}$
%%%END
%%%BLOCK id=238-01 ch=02 sec=力的合成 kind=knowledge alt=none cont=none y=0.04-0.27
\textbf{力的合成}

求合力　正弦定理　余弦定理

\[
F=\sqrt{F_1^2+F_2^2+2F_1F_2\cos\theta}
\]
随$\theta\uparrow$而$\downarrow$　　$\theta=180^\circ$ 最小 $F_1-F_2$　　$\theta=0^\circ$ 最大 $F_1+F_2$

%FIG p238_a 0.25 0.168 0.535 0.24
\notefig[0.3]{figures/p238_a.png}

首尾相连成闭合图形　$F_{\text{合}}=0$
%%%END
%%%BLOCK id=238-02 ch=02 sec=力的分解·动态平衡 kind=method alt=none cont=none y=0.27-0.45
力的分解、画动态图求力　　按实际作用效果正交分解

%FIG p238_b 0.10 0.39 0.205 0.45
\notefig[0.25]{figures/p238_b.png}

$N_{\text{地B}}$不变　　$T\uparrow$
%%%END
%%%BLOCK id=238-03 ch=02 sec=力的分解·动态平衡 kind=method alt=none cont=none y=0.42-0.57
%FIG p238_c 0.09 0.46 0.335 0.562
\notefig[0.3]{figures/p238_c.png}

%FIG p238_d 0.395 0.415 0.70 0.54
\notefig[0.3]{figures/p238_d.png}

$N_{OA}\uparrow$

$N_{OB}$先$\uparrow$后$\downarrow$
%%%END
%%%BLOCK id=238-04 ch=02 sec=共点力平衡·摩擦临界 kind=derivation alt=none cont=none y=0.59-0.70
%FIG p238_e 0.055 0.598 0.225 0.688
\notefig[0.25]{figures/p238_e.png}

\[
2\mu N\cos\theta=mg+2N\sin\theta
\]
\[
N=\frac{mg}{2(\mu\cos\theta-\sin\theta)}
\]
想\unclear{卡住}则 $N>0$　　$\theta\uparrow\ N\uparrow$ % CHECK: "想?则"中间一字被涂改，无法辨认

$\mu>\tan\theta$
%%%END
%%%BLOCK id=238-05 ch=02 sec=共点力平衡·斜面受力分析 kind=note alt=03 cont=none y=0.71-0.80
%FIG p238_f 0.07 0.74 0.255 0.81
\notefig[0.25]{figures/p238_f.png}

加$F$　$f_{\text{斜物}}\uparrow$　$N_{\text{斜物}}\uparrow$　$f_{\text{地斜}}=0$　$N_{\text{地斜}}\uparrow$
%%%END
%%%BLOCK id=238-06 ch=02 sec=共点力平衡·临界极值 kind=problem alt=none cont=none y=0.82-0.98
\begin{problem}
使\unclear{球}重力加$\Delta G$ 还使$C$静止　$\Delta G_{\max}=?$ % CHECK: "球"字手写像"挂"，按语义转录

%FIG p238_g 0.06 0.84 0.235 0.918
\notefig[0.2]{figures/p238_g.png}

\begin{solution}
\begin{align*}
F_N&=\sqrt2\,(G+\Delta G_m) & \Delta G&=\frac{2\mu-1}{1-\mu}\,G\\
F_N\sin45^\circ&=\mu(G+F\cos45^\circ)
\end{align*}
\end{solution}
\end{problem}
%%%END
%%%BLOCK id=240-01 ch=02 sec=共点力平衡·活结与死结 kind=knowledge alt=none cont=none y=0.01-0.15
% 左上角淡铅笔小字标题
\textbf{活结}

%FIG p240_a 0.06 0.03 0.17 0.108
\notefig[0.2]{figures/p240_a.png}

活结两侧拉力相等

%FIG p240_b 0.222 0.046 0.318 0.102
\notefig[0.2]{figures/p240_b.png}

$T\cos\theta=\frac12 mg$　　$T=\dfrac{mg}{2\cos\theta}$　　$2T\cos\theta=mg$　　$d\uparrow\ \cos\alpha\downarrow\ T\uparrow$ % CHECK: 末式中的角手写像 α（前式为 θ），照录

$\sin\theta=\dfrac{d}{L}$

%FIG p240_c 0.555 0.108 0.645 0.15
\notefig[0.2]{figures/p240_c.png}

引体向上　\unclear{竖}手力小
%%%END
%%%BLOCK id=240-02 ch=02 sec=共点力平衡·活结与死结 kind=knowledge alt=none cont=none y=0.16-0.25
%FIG p240_d 0.075 0.158 0.205 0.25
\notefig[0.25]{figures/p240_d.png}

$d$先$\uparrow$后$\downarrow$

$T$先$\uparrow$后$\downarrow$
%%%END
%%%BLOCK id=240-03 ch=02 sec=共点力平衡·活结与死结 kind=knowledge alt=none cont=none y=0.25-0.37
死结　结不沿绳移动，为两根绳

活杆　沿杆

死杆　任意

%FIG p240_e 0.52 0.255 0.68 0.368
\notefig[0.2]{figures/p240_e.png}
%%%END
%%%BLOCK id=240-04 ch=02 sec=共点力平衡·条件与方法 kind=knowledge alt=none cont=none y=0.38-0.46
共点力（延长线交于1点）平衡 $\Leftrightarrow$ 缓慢　匀速、静止、$F_{\text{合}}=0$ % 原稿"1点"写在括号内第二行

方法：二力平衡条件　等效
%%%END
%%%BLOCK id=240-05 ch=02 sec=共点力平衡·条件与方法 kind=method alt=03 cont=none y=0.47-0.57
%FIG p240_f 0.10 0.488 0.245 0.568
\notefig[0.2]{figures/p240_f.png}

平衡后剪断绳子，$F_B=mg$

%FIG p240_g 0.575 0.462 0.835 0.55
\notefig[0.3]{figures/p240_g.png}
% 斜面简图旁手写：α t_1, t_2 α（斜面上方），t_1（块旁）；弯箭头自"验证"处指向左侧 F_B=mg
$\alpha\ \checkmark$　验证
%%%END
%%%BLOCK id=240-06 ch=02 sec=共点力平衡·条件与方法 kind=method alt=none cont=none y=0.57-0.66
三力平衡　力三角形

多力平衡　正交分解

相似，力辅助圆，正弦定理、余弦定理　勾股定理　　$\sin(\pi-\alpha)=\sin\alpha$
%%%END
%%%BLOCK id=240-07 ch=02 sec=共点力平衡·摩擦角 kind=derivation alt=none cont=none y=0.67-0.78
$4\to3$

%FIG p240_h 0.14 0.683 0.44 0.78
\notefig[0.3]{figures/p240_h.png}

\[
F=mg\sin\varphi=\frac{\mu mg}{\sqrt{1+\mu^2}}
\]
%%%END
%%%BLOCK id=241-01 ch=02 sec=页眉·力 kind=note alt=none cont=none y=0.06-0.12
\textbf{力}\quad（页面左上角手写的大字标题）
%%%END
%%%BLOCK id=241-03 ch=02 sec=整体法隔离法·楔形体与球 kind=diagram alt=none cont=none y=0.16-0.34
\begin{problem}
%FIG p241_b 0.13 0.16 0.60 0.255
\notefig[0.47]{figures/p241_b.png}
%FIG p241_c 0.135 0.255 0.23 0.34
\notefig[0.25]{figures/p241_c.png}
\begin{solution}
\[ F\downarrow\qquad F'\downarrow\qquad E\downarrow \] % CHECK: 第三个符号为“E”（可能是 F'' 或 N）
\end{solution}
\end{problem}
%%%END
%%%BLOCK id=241-04 ch=02 sec=静摩擦力·方向讨论 kind=problem alt=none cont=none y=0.38-0.49
\begin{problem}
%FIG p241_d 0.10 0.415 0.25 0.485
\notefig[0.25]{figures/p241_d.png}
\begin{solution}
\[ mg\sin\theta=F\cos\theta\pm f_{12} \]
\[ f_{BA}\text{不定}\ \begin{cases}\text{先}\uparrow\text{后}\downarrow\\ \uparrow\end{cases}\qquad\qquad f_{\text{地}}=F\uparrow \]
\unclear{静摩擦力}永远\unclear{要方}向讨论 % CHECK: 此句在 f_地=F↑ 上方，前半句笔迹潦草（静?止?摩擦…力），整句大意“静摩擦力方向永远要讨论”
\end{solution}
\end{problem}
%%%END
%%%BLOCK id=241-05 ch=02 sec=动态平衡·绳与杆 kind=problem alt=none cont=none y=0.50-0.62
\begin{problem}
%FIG p241_e 0.08 0.50 0.23 0.62
\notefig[0.25]{figures/p241_e.png}
\begin{solution}
\[ 2F\cos\theta=mg\qquad F_1=F_2=mg\quad\text{又绳力不变}\therefore\text{杆力}\uparrow \] % CHECK: “∴”原稿像“心”
\[ A\text{向下}\quad\theta\uparrow\quad F\uparrow \]
\end{solution}
\end{problem}
%%%END
%%%BLOCK id=241-06 ch=02 sec=摩擦力·圆柱放在斜轨上 kind=problem alt=none cont=none y=0.62-0.76
\begin{problem}
%FIG p241_f 0.08 0.606 0.26 0.745
\notefig[0.3]{figures/p241_f.png}
\begin{solution}
\[ 2\mu mg\cos\theta=mg\sin\theta\qquad \mu=2\tan\theta \] % CHECK: 按上式应得 μ=½tanθ，原稿写作“2tanθ”
\[ N_{\text{左}}=N_{\text{右}}=mg\cos\theta\qquad f_{\text{左}}=f_{\text{右}}=\ \tfrac12 mg\sin\theta \] % 等号与 1/2 之间原稿有一块被涂白/遮盖的空白区（可能盖住了字）
\[ \theta\uparrow\quad N\downarrow\quad f\uparrow \]
\end{solution}
\end{problem}
%%%END
%%%BLOCK id=241-07 ch=02 sec=动态平衡·结点与绳OA kind=problem alt=none cont=none y=0.76-0.88
\begin{problem}
%FIG p241_g 0.06 0.765 0.265 0.885
\notefig[0.25]{figures/p241_g.png}
$O$从$A$到$B$ 求$F_{OA}$ % CHECK: 首字为“O”（也可能是①）
\begin{solution}
$OA$在$\angle BOC$\unclear{平}分线上\qquad $F_{OA}=2F\cos\theta$\qquad $\theta\uparrow$\quad $F_{OA}\downarrow$
\end{solution}
\end{problem}
%%%END
%%%BLOCK id=241-08 ch=02 sec=共点力平衡·双绳与质量比 kind=problem alt=none cont=none y=0.88-1.00
\begin{problem}
%FIG p241_h 0.05 0.885 0.205 0.995
\notefig[0.25]{figures/p241_h.png}
\begin{solution}
\begin{align*}
T\sin\alpha&=m_Ag\\
T\sin\beta&=m_Bg
\end{align*}
\[ \frac{m_A}{m_B}=\frac{\sin\alpha}{\sin\beta} \]
\end{solution}
\end{problem}
%%%END
%%%BLOCK id=242-01 ch=02 sec=动态平衡·两绳夹角不变 kind=problem alt=none cont=none y=0.00-0.24
\begin{problem}[1]
%FIG p242_a 0.07 0.00 0.235 0.072
\notefig[0.25]{figures/p242_a.png}
%FIG p242_b 0.30 0.00 0.40 0.05
\notefig[0.25]{figures/p242_b.png}
%FIG p242_c 0.12 0.075 0.29 0.15
\notefig[0.25]{figures/p242_c.png}
\begin{solution}
\[ \begin{cases} F_1\sin\alpha=F_2\sin(\theta-\alpha)\\ F_1\cos\alpha+F_2\cos(\theta-\alpha)=G \end{cases} \]
\[ F_1=G\,\frac{\sin(\theta-\alpha)}{\sin\theta}\qquad F_2=G\,\frac{\sin\alpha}{\sin\theta}\qquad\theta\text{不变} \]
\begin{itemize}
\item $\alpha$从$90^\circ\to0^\circ$
\item $\theta-\alpha$\unclear{在}\ $120^\circ$\ 从$30^\circ\sim120^\circ$ % CHECK: 此处三行斜排：“θ−α ?”“120°”“从30°~120°”，120°应指θ=120°
\item $F_1$先$\uparrow$后$\downarrow$
\item $F_2\downarrow\to0$
\end{itemize}
\end{solution}
\end{problem}
%%%END
%%%BLOCK id=242-05 ch=02 sec=共点力平衡·相似三角形 kind=problem alt=none cont=none y=0.70-0.865
\begin{problem}[5]
%FIG p242_g 0.07 0.70 0.225 0.865
\notefig[0.25]{figures/p242_g.png}
\begin{solution}
\[ \frac{T}{l_1}=\frac{m_1g}{h}\qquad\frac{T}{l_2}=\frac{m_2g}{h} \]
\[ \therefore m_1l_1=m_2l_2 \]
\end{solution}
\end{problem}
%%%END
%%%BLOCK id=242-06 ch=02 sec=共点力平衡·靠墙杆的受力 kind=problem alt=none cont=none y=0.865-1.00
\begin{problem}[6]
%FIG p242_h 0.11 0.865 0.245 0.985
\notefig[0.25]{figures/p242_h.png}
\begin{solution}
法1\ 用杆\unclear{力矩}平衡 % CHECK: “力矩”二字笔迹潦草，仅辨认出大意
\[ \underset{\uparrow}{F_1}\,OA=G\,O\underset{\uparrow}{B} \] % 原稿 F_1 与 B 的下方各有一个向上的小箭头
\[ N_2=G\qquad f=F_1 \]
地对木杆力\unclear{即}$\uparrow$ $\vec N+\vec f$

%FIG p242_i 0.46 0.88 0.585 0.96
\notefig[0.25]{figures/p242_i.png}

法2\ \unclear{竖水平分} % CHECK: 末字不清（也可能是“衡”），大意应为竖直/水平方向分解
%FIG p242_j 0.72 0.915 0.885 0.98
\notefig[0.25]{figures/p242_j.png}
\end{solution}
\end{problem}
%%%END
%%%BLOCK id=246-01 ch=02 sec=共点力平衡·动态平衡 kind=diagram alt=none cont=none y=0.00-0.14
%FIG p246_a 0.02 0.00 0.24 0.145
\notefig[0.25]{figures/p246_a.png}
$A$先$\downarrow$后$\uparrow$，$B\uparrow$\quad $A\uparrow$ % 末尾“A”旁的箭头为弯曲向上的箭头
%%%END
%%%BLOCK id=246-02 ch=02 sec=共点力平衡·受力分析 kind=problem alt=none cont=none y=0.02-0.24
\begin{problem}
现在平衡，求 $\tan\theta$。
%FIG p246_b 0.34 0.03 0.73 0.14
\notefig[0.4]{figures/p246_b.png}
\begin{solution}
\[
\begin{cases}
N\sin37^\circ=F\cos\theta\\
N\cos37^\circ=m_1g+F\cos\theta\\
2F\sin\theta=m_1g
\end{cases}
\qquad \tan\theta=\dfrac{16}{27}
\]
\end{solution}
\end{problem}
%%%END
%%%BLOCK id=247-04 ch=02 sec=受力分析 kind=summary alt=03 cont=none y=0.76-0.83
\textbf{力学两确定}
\begin{itemize}
\item 平衡非平衡
\item 研究对象
\end{itemize}
%%%END
%%%BLOCK id=247-05 ch=02 sec=受力分析 kind=summary alt=none cont=none y=0.86-1.00
\textbf{受力分析}
\begin{itemize}
\item 力三角形（旋转、延长）
\item 非特殊三角形\quad 相似、正余弦定理
\item 多力正交分解
\item 定角定边画力辅助圆
\end{itemize}
%%%END
%%%BLOCK id=249-03 ch=02 sec=共点力平衡·动态平衡 kind=problem alt=none cont=none y=0.69-0.81
%FIG p249_d 0.04 0.685 0.256 0.81
\notefig[0.25]{figures/p249_d.png}
%FIG p249_e 0.44 0.715 0.56 0.795
\notefig[0.25]{figures/p249_e.png}
恒力 $F$ 从 $OM$ 转到 $ON$ \\
$T_{OA}$ 先$\uparrow$后$\downarrow$ \\
$T_{OB}\uparrow$
%%%END
%%%BLOCK id=258-05 ch=02 sec=受力分析·力学两确定 kind=method alt=03 cont=none y=0.55-0.62
\textbf{力学两确定}
\begin{itemize}
\item 状态
  \begin{itemize}
  \item 平衡
  \item 非平衡
  \end{itemize}
\item 研究对象
\end{itemize}
%%%END
%%%BLOCK id=258-06 ch=02 sec=共点力平衡·动态平衡 kind=summary alt=none cont=none y=0.62-0.88
\textbf{动态平衡}
% 原稿：前三项另有一小括号相连，七项再由一个大括号统领
\begin{itemize}
\item 力三角形、平移、延长\ 旋转
\item 正弦、余弦定理、\unclear{余力}勾股定理\ 三角函数
\item 非特殊三角形\ 用相似
\item 定角画圆
\item 晾衣\unclear{架}模型\ $+$\ 转动 % 「架」处原稿字迹有叠写
\item 正交分解
\item 摩擦角、全反力
\end{itemize}
%%%END
%%%BLOCK id=259-01 ch=02 sec=受力分析·摩擦力 kind=diagram alt=none cont=none y=0.00-0.15
\textbf{人撑在两面垂直的墙}
%FIG p259_a 0.07 0.00 0.25 0.15
\notefig[0.3]{figures/p259_a.png}
%%%END
%%%BLOCK id=259-06 ch=02 sec=受力分析·多物体 kind=method alt=03 cont=none y=0.76-0.83
多物体复杂受力分析\quad 质量分块、受力均摊，分析节点（重力可能不是 $mg$，是\unclear{等效}结果 % 此行写到页边被截断，括号未闭合
\par 几何、三角函数关系 % 原稿写在上一行「分析节点」上方
%%%END
%%%BLOCK id=267-03 ch=02 sec=共点力平衡·摩擦角法 kind=method alt=none cont=none y=0.31-0.48
\textbf{使用摩擦角}
%FIG p267_b 0.02 0.305 0.33 0.485
\notefig[0.45]{figures/p267_b.png}
$f\downarrow$\quad $F$ 先 $\downarrow$ 后 $\uparrow$.
%%%END
%%%BLOCK id=269-05 ch=02 sec=力的合成·验证平行四边形定则 kind=note exp=1 alt=none cont=none y=0.90-0.93
\unclear{验}力的平行四边形定则时\ 不在平行四边形中的是理想的
%%%END
%%%BLOCK id=271-04 ch=02 sec=弹力·等效劲度系数 kind=knowledge alt=06 cont=none y=0.47-0.61
%FIG p271_c 0.555 0.485 0.725 0.61
\notefig[0.25]{figures/p271_c.png}
求系统等效劲度系数。图中标注：$2x\downarrow$，$2kx$；$\downarrow4kx$；$x\downarrow$。
\[
K_{\text{等效}}=4K
\]
%%%END
%%%BLOCK id=280-02 ch=02 sec=全反力·匀速运动 kind=method alt=03 cont=none y=0.24-0.46
%FIG p280_b 0.09 0.24 0.36 0.385
\notefig[0.3]{figures/p280_b.png}
$\mu=\dfrac{\sqrt{3}}{3}$\quad 匀速运动

%FIG p280_c 0.38 0.28 0.65 0.45
\notefig[0.3]{figures/p280_c.png}
图中标注：$mg$，$R$ 全\unclear{反}力

$F\uparrow\qquad N\uparrow\qquad f\uparrow$
%%%END
%%%BLOCK id=280-04 ch=02 sec=三力汇交·杆与球形碗 kind=problem alt=none cont=none y=0.69-0.87
支持力$N$与接触面垂直

%FIG p280_e 0.18 0.725 0.52 0.865
\notefig[0.4]{figures/p280_e.png}
\unclear{多}力交汇% CHECK: 首字疑为“三”（三力汇交）
\[
2R\cos\theta-2R\sin\theta\tan\theta=\frac{L}{2}\qquad \therefore L=\frac{4R\cos2\theta}{\cos\theta}
\]
%%%END
%%%BLOCK id=285-08 ch=02 sec=叠放物体·静摩擦与做功 kind=derivation alt=06 cont=none y=0.36-0.49
%FIG p285_b 0.15 0.365 0.335 0.45
\notefig[0.25]{figures/p285_b.png}

\begin{gather*}
2F_m\cos60^\circ=mg\qquad F_m=mg\\
f_m=F_m\cos30^\circ=\frac{\sqrt{3}}{2}mg\\
\text{B对地压力}\ F_N=m_Bg+\tfrac12 m_Cg=mg\\
f_{\text{地B}}=f_m\qquad \mu=\frac{\sqrt{3}}{2}
\end{gather*}
% CHECK: F_N 一行中 m 的下标手写不清（读作 B、C）
% 以下三式写在上述各式右侧
\begin{gather*}
W-fx+mgh=0\\
fx=2(\sqrt{3}-1)\mu mgR\\
W=(2\mu-1)(\unclear{\sqrt{3}}-1)mgR
\end{gather*}
% 第二式中 x 右下有一个实心小点（可能是下标0）；第三式根号内手写像 2，按上下文读作 3
%%%END
%%%BLOCK id=286-02 ch=02 sec=叠放物体·受力分析 kind=note alt=none cont=none y=0.63-0.71
%FIG p286_b 0.09 0.635 0.205 0.71
\notefig[0.25]{figures/p286_b.png}

$F_{2\text{地}}=2.5G$\quad 两边不算
%%%END
